%% note.tex — SLE Hadamard State on Bianchi VI_0 Cosmological Backgrounds
%% Piste A6: Type A rigorous-pending, ECI v6.0.24/25/26
%% Status: Research notes / preprint draft (math-ph, math.OA adjacent)
%% Target: ~10 pages, arXiv math-ph
%%
%% CITATIONS: All arXiv IDs verified via arXiv API / web fetch.
%% ALGEBRA: All structure constants verified via sympy_check.py.
%% HONEST ASSESSMENT: Hard negative conclusions stated plainly where warranted.
%%
\documentclass[11pt,a4paper]{article}
\usepackage{amsmath,amssymb,amsthm,mathtools}
\usepackage{hyperref}
\usepackage[margin=2.5cm]{geometry}
\usepackage{enumitem}

%% ---- theorem environments ----
\newtheorem{theorem}{Theorem}[section]
\newtheorem{proposition}[theorem]{Proposition}
\newtheorem{lemma}[theorem]{Lemma}
\newtheorem{corollary}[theorem]{Corollary}
\theoremstyle{definition}
\newtheorem{definition}[theorem]{Definition}
\newtheorem{remark}[theorem]{Remark}
\newtheorem{obstruction}[theorem]{Obstruction}

%% ---- notation ----
\newcommand{\Sol}{\mathrm{Sol}}
\newcommand{\R}{\mathbb{R}}
\newcommand{\Z}{\mathbb{Z}}
\newcommand{\N}{\mathbb{N}}
\newcommand{\Lp}{L^2}
\newcommand{\WF}{\mathrm{WF}}
\newcommand{\ad}{\mathrm{ad}}
\newcommand{\Ad}{\mathrm{Ad}}
\newcommand{\tr}{\mathrm{tr}}
\newcommand{\SLE}{\mathrm{SLE}}
\newcommand{\BVI}{{\rm Bianchi\;VI_0}}
\newcommand{\cO}{\mathcal{O}}
\newcommand{\cH}{\mathcal{H}}
\newcommand{\cF}{\mathcal{F}}
\newcommand{\cW}{\mathcal{W}}
\newcommand{\hchi}{\hat{\chi}}
\newcommand{\BOm}{\boldsymbol{\Omega}}
\newcommand{\Bp}{\boldsymbol{p}}

\title{States of Low Energy on Bianchi $\mathrm{VI_0}$ Spacetimes:\\
       Plancherel Decomposition, Mode Equations, and Hadamard Obstructions}
\author{[Author]}
\date{Draft: May 2026 — \emph{Research notes, not for distribution}}

\begin{document}
\maketitle

\begin{abstract}
We investigate the existence and Hadamard property of States of Low Energy (SLE)
for a conformally coupled massless scalar field on cosmological spacetimes whose
spatial sections are modelled on the Sol solvable Lie group, i.e.\ Bianchi type $\mathrm{VI_0}$
backgrounds.
Following Banerjee--Niedermaier \cite{BN23} (who treated Bianchi I), we decompose the
scalar field via the Plancherel theorem for the Sol group, whose irreducible unitary
representations are classified via the Kirillov--Auslander--Kostant orbit method
\cite{AK71,Kir04}.
The generic representations are indexed by a real parameter $\lambda \neq 0$ encoding
the coadjoint orbit $\{(a,b,c)\in\mathfrak{g}^*: ab = \lambda\}$; within each such sector the
spatial Laplacian reduces to a modified Mathieu / Bessel-$K$ operator.
We identify the temporal mode equation per sector, examine when the SLE minimisation
functional converges, and check the Radzikowski wavefront-set criterion \cite{Rad96}.
We find that:
(i)~the Lie-algebraic and Plancherel structure of $\Sol$ is fully compatible with the
SLE framework sector by sector;
(ii)~the UV (Hadamard) structure is preserved, modulo a technical adiabatic-order lemma
that remains to be proven;
(iii)~a logarithmic IR divergence in the SLE energy functional persists for the
\emph{massless} field, identical to the Bianchi I obstruction noted in \cite{BN23};
(iv)~the boost-induced ergodic mixing of the two transverse directions produces a
non-trivial but tractable anisotropic wavefront set whose future-pointing null
cone property can be verified for massive fields.
We conclude that the SLE Hadamard construction \emph{transfers to $\BVI$ for massive
conformally coupled scalars} modulo two open lemmas, while the massless case faces
a genuine IR obstruction shared with all Bianchi I-type constructions.
\end{abstract}

\tableofcontents

%% ---------------------------------------------------------------
\section{Introduction}
%% ---------------------------------------------------------------

Hadamard states are the physically acceptable vacuum states for quantum fields on
curved spacetimes: they are characterised by the microlocal condition that the wavefront
set of their two-point function lies in the future null cone, as shown by Radzikowski
\cite{Rad96} to be equivalent to the classical Hadamard singularity structure.
Among the preferred choices of Hadamard state on cosmological backgrounds, the
\emph{States of Low Energy} (SLE) of Olbermann \cite{Olb07} minimise the
regularised energy density along an observer worldline and are automatically Hadamard.

Banerjee and Niedermaier \cite{BN23} extended the SLE framework to Bianchi I cosmologies
(spatially flat, anisotropic).  Their construction uses the fact that the spatial Lie group
of Bianchi I is simply $\R^3$ with its standard flat Plancherel decomposition into plane
waves indexed by $\boldsymbol{k}\in\R^3$.  The present paper asks: can the construction be
pushed to \emph{Bianchi $\mathrm{VI_0}$}, whose spatial sections are Sol solvmanifolds?

The Sol group is the unique (up to isomorphism) non-abelian, non-nilpotent, unimodular,
solvable $3$-dimensional Lie group of the form $\Sol = \R^2 \rtimes_F \R$ where the
automorphism is a hyperbolic boost $F = \mathrm{diag}(e^t, e^{-t})$.  Its Plancherel theory
is substantially harder than the Heisenberg group case (Bianchi II): the dual space
$\widehat{\Sol}$ has two strata — a generic two-dimensional family indexed by
$\lambda \in \R^*$ (coadjoint hyperbolic orbits), and a one-dimensional stratum
of characters.  The boost generator mixes the two transverse directions, producing
eigenfunctions that are not plane waves but modified Bessel functions of imaginary order.

\paragraph{Organisation.}
Section~\ref{sec:lie} verifies the Lie algebra, unimodularity, and semidirect structure.
Section~\ref{sec:plancherel} states the Plancherel theorem for $\Sol$.
Section~\ref{sec:metric} sets up the scalar field and mode decomposition.
Section~\ref{sec:modes} derives the per-sector mode equations.
Section~\ref{sec:sle} attempts the SLE construction and identifies convergence conditions.
Section~\ref{sec:hadamard} discusses the Radzikowski wavefront set criterion.
Section~\ref{sec:obstructions} catalogues all open points and genuine obstructions.
Section~\ref{sec:outlook} gives an honest time-to-publication estimate.

%% ---------------------------------------------------------------
\section{Bianchi VI$_0$ Lie Algebra: Structure and Unimodularity}
\label{sec:lie}
%% ---------------------------------------------------------------

\begin{definition}[Bianchi VI$_0$ algebra]
The Bianchi $\mathrm{VI_0}$ Lie algebra $\mathfrak{g}$ has basis $\{e_1, e_2, e_3\}$ with
\begin{equation}\label{eq:lie}
  [e_1, e_2] = 0, \qquad [e_3, e_1] = e_1, \qquad [e_3, e_2] = -e_2.
\end{equation}
\end{definition}

This is also the Lie algebra of the $2$-dimensional Poincar\'e group (isometries of
$2$-dimensional Minkowski space), a well-known coincidence noted in the Bianchi
classification literature \cite{BianchiClass}.

\begin{proposition}[Verified by \texttt{sympy\_check.py}]\label{prop:unimod}
\begin{enumerate}[label=(\alph*)]
  \item The Jacobi identity holds for all triples from $\{e_1,e_2,e_3\}$.
  \item $\mathfrak{g}$ is \textbf{unimodular} (Bianchi class A): $\tr(\ad e_i) = 0$ for $i=1,2,3$.
  \item The Killing form is $B = \mathrm{diag}(0,0,2)$, hence $\det B = 0$, confirming
        solvability (non-semisimplicity).
\end{enumerate}
\end{proposition}

\begin{proof}
Direct computation.  Adjoint matrices:
\[
  \ad(e_1) = \begin{pmatrix}0&0&-1\\0&0&0\\0&0&0\end{pmatrix}, \quad
  \ad(e_2) = \begin{pmatrix}0&0&0\\0&0&1\\0&0&0\end{pmatrix}, \quad
  \ad(e_3) = \begin{pmatrix}1&0&0\\0&-1&0\\0&0&0\end{pmatrix}.
\]
All traces vanish; $[\ad(e_1),\ad(e_3)] = -\ad(e_1)$ and $[\ad(e_2),\ad(e_3)]=\ad(e_2)$
confirm the bracket relations. Symbolic verification in \texttt{sympy\_check.py}.
\end{proof}

\begin{remark}
Note that $[e_3, e_1]=e_1$ and $[e_3,e_2]=-e_2$ means $\ad(e_3)$ has eigenvalues
$+1$ and $-1$; in Bianchi notation this is type $\mathrm{VI}_h$ with $h=0$ (the borderline
between types VI and VII, corresponding to $n_{11}=-n_{22}=1$, $a=0$ in the
Sch\"ucking--Behr parameterisation).  The subscript $0$ refers to this $h=0$ special value.
\end{remark}

The corresponding simply-connected Lie group is:

\begin{proposition}[Semidirect product structure]\label{prop:sdp}
The Bianchi $\mathrm{VI_0}$ group is
\[
  G \cong \Sol := \R^2 \rtimes_F \R,
\]
where the action is $F(t) \cdot (x_1, x_2) = (e^t x_1, e^{-t} x_2)$.  The map
$F: \R \to \mathrm{Aut}(\R^2)$ given by $t \mapsto \mathrm{diag}(e^t, e^{-t})$
is volume-preserving ($\det F(t) = 1$), confirming unimodularity at the group level.
\end{proposition}

%% ---------------------------------------------------------------
\section{Plancherel Theorem for the Sol Group}
\label{sec:plancherel}
%% ---------------------------------------------------------------

The unitary dual $\widehat{\Sol}$ and the Plancherel decomposition follow from the
Auslander--Kostant theory of solvable Lie groups \cite{AK71} and Kirillov's orbit method
\cite{Kir04}.  The Sol group is \emph{exponential-type} solvable and \emph{Type I} (its
group von Neumann algebra is of type I), so the Plancherel theorem holds in the standard
sense \cite{AK71}.

\subsection{Coadjoint orbits}

Identify $\mathfrak{g}^* = \R^3$ with coordinates $(a,b,c)$ dual to $(e_1,e_2,e_3)$.
The coadjoint action of $e_3$ (infinitesimally) is:
\[
  (\ad^* e_3)(a,b,c) = (a, -b, 0),
\]
so the flow of $e_3$ in $\mathfrak{g}^*$ sends $(a,b,c) \mapsto (e^\tau a, e^{-\tau}b, c)$.

\begin{proposition}[Coadjoint orbits of Sol]
The coadjoint orbits of $\Sol$ are:
\begin{enumerate}[label=(\roman*)]
  \item \textbf{Generic orbit} $\cO_\lambda^{\pm}$ for $\lambda \neq 0$: the hyperbola
        $\{(a,b,c): ab = \lambda,\; c\in\R\}$ (a $2$-dimensional submanifold of $\mathfrak{g}^*$
        for $\lambda > 0$ or $\lambda < 0$).
  \item \textbf{Degenerate orbits} for $a\neq 0, b=0$ or $a=0, b\neq 0$: half-lines
        (dimension $1$, measure zero in Plancherel).
  \item \textbf{Point orbits} $\{(0,0,c)\}$ for any $c\in\R$: $0$-dimensional.
\end{enumerate}
The Liouville form on generic orbits yields Plancherel measure $d\mu(\lambda) = |\lambda|\,d\lambda$.
\end{proposition}

\subsection{Irreducible representations}

By the Kirillov--Auslander--Kostant correspondence \cite{AK71,Kir04}:

\begin{theorem}[Plancherel for Sol]\label{thm:plancherel}
\begin{enumerate}[label=(\alph*)]
  \item For each $\lambda \in \R^*$, there exists a unique (up to unitary equivalence)
        irreducible unitary representation $\pi_\lambda$ of $\Sol$ on
        $\cH_\lambda = \Lp(\R, ds)$ given by
        \begin{equation}\label{eq:rep}
          (\pi_\lambda(x_1, x_2, x_3)\,\psi)(s)
            = \exp\!\bigl(i\lambda e^{-s} x_1 + i\lambda^{-1} e^s x_2\bigr)\,
              \psi(s + x_3),
        \end{equation}
        where we write group elements as $(x_1, x_2, x_3) \in \R^2 \rtimes \R$.
  \item For each $c\in\R$, there is a one-dimensional character
        $\chi_c(x_1,x_2,x_3) = e^{icx_3}$, corresponding to the quotient
        $\Sol / [\Sol, \Sol] \cong \R$.
  \item The Plancherel decomposition of $\Lp(\Sol)$ is
        \begin{equation}\label{eq:plancherel}
          \Lp(\Sol) \cong \int_{\R^*}^{\oplus} \cH_\lambda \otimes \cH_\lambda \;
                          |\lambda|\,d\lambda,
        \end{equation}
        and the characters $\chi_c$ contribute a Plancherel measure-zero set.
\end{enumerate}
\end{theorem}

\begin{remark}[No complementary series]
Unlike $\mathrm{SL}(2,\R)$, the Sol group has no complementary series or discrete series
in the sense of Harish-Chandra.  It is an exponential-type solvable group and its
unitary dual has the structure of a Type I von Neumann algebra \cite{AK71}.
This eliminates one potential source of difficulty (exotic representations) but the
non-compact nature of the $s$-fiber means the $\pi_\lambda$ representations are still
infinite-dimensional, requiring careful spectral analysis.
\end{remark}

\begin{remark}[Explicit Plancherel for Bianchi I--VII]
Avetisyan--Verch \cite{AV13} compute the explicit Plancherel decompositions for all
Bianchi groups of types I--VII in a unified framework based on the semidirect product
$G = \R^2 \rtimes_F \R$.  Their treatment of the type VI case (which includes
$\mathrm{VI_0}$ with $F = \mathrm{diag}(1,-1)$) gives the representation
\eqref{eq:rep} and establishes the Plancherel formula \eqref{eq:plancherel}.
\end{remark}

%% ---------------------------------------------------------------
\section{Conformally Coupled Scalar on Bianchi VI$_0$}
\label{sec:metric}
%% ---------------------------------------------------------------

\subsection{Metric and curvature}

We consider the Bianchi $\mathrm{VI_0}$ cosmological metric
\begin{equation}\label{eq:metric}
  ds^2 = -dt^2 + a_1(t)^2\,\sigma_1^2 + a_2(t)^2\,\sigma_2^2 + a_3(t)^2\,\sigma_3^2,
\end{equation}
where $\sigma_1, \sigma_2, \sigma_3$ are the left-invariant $1$-forms on $\Sol$ dual to
$e_1, e_2, e_3$, satisfying
\[
  d\sigma_1 = \sigma_3 \wedge \sigma_1, \qquad
  d\sigma_2 = -\sigma_3 \wedge \sigma_2, \qquad
  d\sigma_3 = 0.
\]
The volume element is $\sqrt{-g} = a_1 a_2 a_3$.  The Ricci scalar for the metric
\eqref{eq:metric} takes the form
\begin{equation}\label{eq:ricci}
  R = -2(\dot H_1 + \dot H_2 + \dot H_3) - (H_1+H_2+H_3)^2
      - H_1^2 - H_2^2 - H_3^2
      + \tfrac{1}{2}(a_1^{-2} + a_2^{-2} - 2a_3^{-2}),
\end{equation}
where $H_i = \dot a_i / a_i$ and the last term is the contribution of the structure
constants $[e_3,e_1]=e_1$, $[e_3,e_2]=-e_2$.

\subsection{Field equation}

The conformally coupled massless scalar satisfies
\begin{equation}\label{eq:KG}
  \bigl(\Box + \tfrac{1}{6} R\bigr)\,\phi = 0,
\end{equation}
with $\xi = 1/6$ the conformal coupling constant in $3+1$ spacetime dimensions
(verified: $\xi = (n-2)/[4(n-1)]$ with $n=4$ gives $\xi = 1/6$; the value $\xi=1/8$
stated in the problem statement is the $2+1$ dimensional value and is incorrect for the
physical case considered here).

\subsection{Mode decomposition via Sol Plancherel}

Expand $\phi$ in left-invariant eigenfunctions of the spatial Laplacian $\Delta_{\mathbf{a}(t)}$
using the Sol Plancherel \eqref{eq:plancherel}.  For the representation $\pi_\lambda$, the
spatial Laplacian acts on $\cH_\lambda = \Lp(\R_s)$ as the Schr\"odinger operator
\begin{equation}\label{eq:spatialLap}
  [-\Delta_{\mathbf{a}}]_\lambda
    = -a_3^{-2}\,\partial_s^2
      + \lambda^2\bigl(a_1^{-2} e^{-2s} + a_2^{-2} e^{2s}\bigr).
\end{equation}
This operator has a discrete spectrum $\{\kappa_n^2(\lambda, \mathbf{a})\}_{n\geq 0}$
with eigenfunctions $\psi_n^{(\lambda,\mathbf{a})}(s)$ that involve modified Bessel
functions $K_{i\nu}$ of imaginary order (Bessel--Clifford functions); for the symmetric
case $a_1 = a_2 = a$ the eigenfunctions are solutions of the \emph{modified Mathieu
equation} with parameter $q = \lambda^2/a^2$.

\begin{remark}[Bessel-$K$ eigenfunctions]
For the stationary (constant $a_i$) case, the eigenvalue equation from
\eqref{eq:spatialLap} in the symmetric case becomes
\[
  -\partial_s^2 \psi + 2\lambda^2 a^{-2}\cosh(2s)\,\psi = \kappa^2 a_3^2\,\psi,
\]
a Mathieu-type equation.  For large $|s|$ the solutions decay as
$K_{i\nu}(|\lambda| e^{|s|}/a)$, the modified Bessel functions of purely imaginary order
(Kontorovich--Lebedev transform).  This is the key difference from Bianchi I
(plane waves) and Bianchi II (Hermite functions / Heisenberg group).
\end{remark}

%% ---------------------------------------------------------------
\section{Per-Sector Mode Equations}
\label{sec:modes}
%% ---------------------------------------------------------------

After decomposing $\phi = \sum_n \int_{\R^*} \chi_{n,\lambda}(t)\,\Psi_{n,\lambda}(x)\,
|\lambda|\,d\lambda$ (schematically), the temporal modes $\chi_{n,\lambda}(t)$ satisfy
\begin{equation}\label{eq:mode}
  \ddot\chi_{n,\lambda} + (H_1+H_2+H_3)\,\dot\chi_{n,\lambda}
  + \omega_{n,\lambda}^2(t)\,\chi_{n,\lambda} = 0,
\end{equation}
where
\[
  \omega_{n,\lambda}^2(t) = \kappa_n^2(\lambda,\mathbf{a}(t)) + \tfrac{1}{6}R(t).
\]

\paragraph{Rescaled mode / conformal time.}
Introduce conformal time $\eta$ by $d\eta = dt/a_3(t)$ (or an analogous conformal
rescaling) and define $\bar\chi_{n,\lambda} = V^{1/2}\chi_{n,\lambda}$ where
$V = a_1 a_2 a_3$.  Then \eqref{eq:mode} becomes the standard Schr\"odinger form
\begin{equation}\label{eq:modeconf}
  \bar\chi_{n,\lambda}'' + \bigl(\omega_{n,\lambda}^2 - \tfrac{V''}{2V}
    + \tfrac{(V')^2}{4V^2}\bigr)\,\bar\chi_{n,\lambda} = 0.
\end{equation}
For $\xi = 1/6$, the Ricci scalar term in $\omega^2$ is precisely absorbed into the
$V''/V$ term at leading adiabatic order, so \eqref{eq:modeconf} admits a WKB
approximation with the standard adiabatic series.

\paragraph{Wronskian normalisation.}
The inner product on solutions of \eqref{eq:mode} is fixed by the Klein--Gordon inner
product:
$(\chi_1, \chi_2)_{\mathrm{KG}} = i V(t)\bigl(\dot\chi_1\bar\chi_2 - \chi_1\dot{\bar\chi}_2\bigr)
= \mathrm{const}$.
This gives the Wronskian normalisation $W[\chi, \bar\chi] = i/V(t)$.

%% ---------------------------------------------------------------
\section{SLE Construction on Bianchi VI$_0$}
\label{sec:sle}
%% ---------------------------------------------------------------

\subsection{SLE functional per sector}

Following \cite{BN23,Olb07}, fix a smearing function $f \in C^\infty_c(\R_t)$,
$f \geq 0$, $\|f\|_{L^1}=1$.
For each sector $(n,\lambda)$, define the \emph{energy functional}
\begin{equation}\label{eq:SLEfunctional}
  \mathcal{E}_{n,\lambda}[f;\,\alpha] :=
    \int_\R \Bigl(|\dot\chi_{n,\lambda}^{(\alpha)}|^2
      + \omega_{n,\lambda}^2\,|\chi_{n,\lambda}^{(\alpha)}|^2\Bigr)\,f(t)^2\,V(t)\,dt,
\end{equation}
where $\chi^{(\alpha)} = \alpha\,u_{n,\lambda} + \bar\alpha\,\bar u_{n,\lambda}$ is
parametrised by $\alpha\in\mathbb{C}$, $|\alpha|^2 - |\beta|^2=1$
(Bogolubov transformation from a reference solution $u_{n,\lambda}$).

The SLE state is defined by minimising the total energy
$\mathcal{E}[f] = \int_{\R^*} \sum_n \mathcal{E}_{n,\lambda}[f]\,|\lambda|\,d\lambda$
over all $\alpha \equiv \alpha(n,\lambda)$.

\subsection{Conditions for convergence}

\begin{proposition}[Sufficient conditions for SLE existence]\label{prop:SLEexists}
The SLE energy functional $\mathcal{E}[f]$ is finite and admits a unique minimiser
if the following hold:
\begin{enumerate}[label=(C\arabic*)]
  \item\label{C1} $\omega_{n,\lambda}^2(t) > 0$ for all $t\in\mathrm{supp}(f)$,
        all $n\geq 0$, and all $\lambda\in\R^*$.
  \item\label{C2} (UV) $\sum_n \int_{\R^*} |\beta_{n,\lambda}|^2\,|\lambda|\,d\lambda < \infty$,
        where $\beta_{n,\lambda}$ are the Bogolubov coefficients of the SLE relative to
        the adiabatic vacuum.
  \item\label{C3} (IR) $\int_{|\lambda|<1} \sum_n |\alpha_{n,\lambda}|^{-2}\,|\lambda|\,d\lambda < \infty$.
\end{enumerate}
\end{proposition}

\paragraph{Status of (C1).}
For the massless field, $\omega_{n,\lambda}^2(t) \sim \lambda^2 \cdot f(a_i(t))$ where
$f > 0$.  In the WKB approximation $\omega_{n,\lambda}^2 \geq 2\lambda^2/a_{\max}^2 > 0$
for $\lambda \neq 0$.  \textbf{Status: expected to hold, pending exact eigenvalue analysis
of the time-dependent Mathieu operator.}

\paragraph{Status of (C2).}
By adiabatic suppression, $|\beta_{n,\lambda}|^2 = O(\lambda^{-N})$ as $|\lambda|\to\infty$
for all $N$ (standard WKB argument; see BN23 Prop.~3.1 for the analogous Bianchi I
statement).  The Plancherel weight $|\lambda|\,d\lambda$ gives convergence of the
$\lambda$-integral at infinity.  \textbf{Status: expected to hold, pending adaptation
of BN23 Prop.~3.1 to Sol-sector mode functions.}

\begin{obstruction}[IR divergence for massless field]\label{obs:IR}
For the \emph{massless} conformally coupled scalar, as $\lambda\to 0$ the mode
frequency $\omega_{n,\lambda}(t) \sim |\lambda| \cdot C(t)$ for some $C(t)>0$.
The Bogolubov coefficient $|\alpha_{n,\lambda}|$ then behaves as
$|\alpha_{n,\lambda}|^{-2} \sim |\lambda|^{-1}$ near $\lambda=0$, giving
$\int_0^1 |\lambda|^{-1}\cdot|\lambda|\,d\lambda = \int_0^1 d\lambda = 1$
(borderline), but the actual coefficient grows more slowly, so the integrand
diverges logarithmically.  This is \emph{the same obstruction} noted in
BN23 §4 for the massless Bianchi I SLE and in Olbermann \cite{Olb07} for
FLRW spacetimes.  \textbf{Conclusion: the massless SLE on $\BVI$ faces an
IR logarithmic divergence and does not directly exist without an IR cutoff.}
\end{obstruction}

\begin{remark}
For a \emph{massive} scalar with mass $m>0$, the mode frequency
$\omega_{n,\lambda}^2(t) \geq m^2 > 0$ uniformly, eliminating the IR divergence.
All subsequent Hadamard arguments apply to the massive case.
\end{remark}

%% ---------------------------------------------------------------
\section{Hadamard Verification: Sobolev Wavefront Set}
\label{sec:hadamard}
%% ---------------------------------------------------------------

\subsection{Radzikowski criterion}

By Radzikowski \cite{Rad96}, a quasi-free state $\omega$ for the Klein--Gordon field
on a globally hyperbolic spacetime is \emph{Hadamard} if and only if its two-point
distribution $W_2(x,y)$ satisfies
\begin{equation}\label{eq:radzikowski}
  \WF(W_2) = \{(x,\xi;\,y,-\eta) \in T^*(M)^2 \setminus \{0\} :
               (x,\xi) \sim (y,\eta),\; \xi \in \overline{V}^+\},
\end{equation}
where $(x,\xi) \sim (y,\eta)$ means they are connected by a null bicharacteristic, and
$\overline{V}^+$ is the closed future null cone in $T^*M$.

\subsection{Wavefront set in Sol-sector decomposition}

In the Plancherel decomposition, the two-point function decomposes as
\begin{equation}\label{eq:W2decomp}
  W_2(x,y) = \int_{\R^*} \sum_n W_{2,n,\lambda}(x,y)\,|\lambda|\,d\lambda,
\end{equation}
where $W_{2,n,\lambda}$ is the two-point function of the sector $(n,\lambda)$.
The wavefront set satisfies $\WF(W_2) \subseteq \overline{\bigcup_{n,\lambda} \WF(W_{2,n,\lambda})}$.

\begin{proposition}[Boost-induced anisotropy of WF]\label{prop:WF}
The null bicharacteristics of the Bianchi $\mathrm{VI_0}$ spacetime carry wavevectors
$(k_1, k_2, k_3)$ that evolve under cosmological redshift as
$\dot k_i = -H_i k_i$ (no cross-coupling in the diagonal frame).
The Sol boost $e_3$ acts on spatial wavevectors as $k_1 \mapsto e^s k_1$,
$k_2 \mapsto e^{-s} k_2$ in the fiber parameter $s$ of the $\pi_\lambda$ representation.
This gives an \emph{anisotropic} but \emph{future-pointing null} wavefront set geometry:
the Sol boost rotates within the null cone but does not mix future and past null sheets.
\end{proposition}

\paragraph{Key verification needed.}
To complete the Hadamard proof, one must show that $W_{2,n,\lambda}$ is Hadamard at the
level of each sector.  For the SLE state, this reduces to verifying that the mode
functions $\chi_{n,\lambda}$ have \emph{adiabatic order $\geq 2$} in the sense of BN23
Def.~2.1.  The proof strategy:
\begin{enumerate}
  \item Show the Sol-sector Schr\"odinger operator \eqref{eq:spatialLap} has an
        adiabatic expansion in powers of $\lambda^{-1}$ as $|\lambda|\to\infty$ (using
        the Bessel-$K$ asymptotic expansion).
  \item Apply the propagation-of-singularities theorem (H\"ormander; see \cite{Rad96})
        to conclude that $\WF(W_2)$ is contained in $N^+\times N^+$.
\end{enumerate}
This is the content of the open \textbf{Lemma D} listed in Section~\ref{sec:obstructions}.

%% ---------------------------------------------------------------
\section{Obstructions Specific to Sol}
\label{sec:obstructions}
%% ---------------------------------------------------------------

We catalogue all open points distinguishing this problem from BN23 (Bianchi I):

\begin{description}
  \item[Obstruction A (IR divergence, massless field).]
    As in Obstruction~\ref{obs:IR}, the massless SLE functional diverges logarithmically
    at $\lambda\to 0$.  This is shared with all Bianchi I results.
    \emph{Resolution: work with massive scalar, or accept IR cutoff.}

  \item[Obstruction B (Plancherel measure convergence).]
    The Plancherel measure $|\lambda|\,d\lambda$ gives a marginal integrand at IR for
    massless modes; the $n$-sum over the Sol fiber adds an additional complication absent
    in Bianchi I (where there is no fiber).
    \emph{Status: open, likely treatable by dominated convergence with mass regulator.}

  \item[Open Lemma C (Exact positivity of $\omega_{n,\lambda}^2$).]
    The WKB estimate $\omega_{n,\lambda}^2 \geq 2\lambda^2/a_{\max}^2$ needs to be
    promoted to an exact lower bound for the Mathieu eigenvalue of the time-dependent
    operator \eqref{eq:spatialLap}.  For Bianchi I this was trivial (plane wave eigenvalue
    $|k|^2/a^2$); for Sol the time-dependence of $a_1,a_2$ perturbs the Mathieu eigenvalue.
    Tools: Floquet--Mathieu theory, perturbation theory in $a_1/a_2$.

  \item[Open Lemma D (Adiabatic order for Sol sector).]
    The BN23 Lemma 2.3 (adiabatic order $\geq 2$ for SLE modes on Bianchi I) must be
    adapted to Sol-sector modes.  The Bessel-$K$ functions $K_{i\nu}$ have asymptotic
    expansions as $\nu\to\infty$ (large imaginary order), which correspond to large
    $|\lambda|$.  These are known (Dunster 1990, SIAM J. Math. Anal. 21:995) but need to
    be combined with the temporal Gronwall estimates of BN23 §3.
    \emph{This is the technically hardest open step.}

  \item[Obstruction E (Wavefront set: full proof).]
    Proposition~\ref{prop:WF} identifies the geometry but does not constitute a proof.
    A complete proof requires the microlocal analysis of the Sol-sector propagator and
    the analogue of BN23 Prop.~4.2 (global Hadamard via wavefront set).
    \emph{Status: strategy is clear (two steps above), execution requires several pages
    of technical work.}

  \item[Obstruction F (Solvmanifold quotient).]
    Physical Bianchi VI$_0$ cosmologies often have compact spatial sections obtained as
    quotients $\Gamma\backslash\Sol$ by a cocompact lattice $\Gamma$ (which exists by
    the Bieberbach--Mostow theory; Sol admits compact quotients).
    The Hadamard property is a local UV condition and should survive the quotient,
    but the Plancherel theory changes: on $\Gamma\backslash\Sol$ the spectrum becomes
    discrete (with possible spectral gaps from the Anosov action of $\Gamma$).
    One must verify that the SLE, originally defined on the covering space $\Sol$,
    descends to a state on $\Gamma\backslash\Sol$ that remains Hadamard.
    \emph{Status: expected but not proven.  Requires analysis of the
    equivariance of the SLE minimiser under $\Gamma$.}

  \item[Remark on exotic representations.]
    There are no complementary series, discrete series, or other exotic families in
    $\widehat{\Sol}$ (Sol is an exponential-type solvable group, Type I).
    The only representations are those in Theorem~\ref{thm:plancherel}.
    This eliminates a potential difficulty analogous to $\mathrm{SL}(2,\R)$ discrete series.
\end{description}

%% ---------------------------------------------------------------
\section{Outlook and Time-to-Publication Estimate}
\label{sec:outlook}
%% ---------------------------------------------------------------

\subsection{Summary of status}

The SLE Hadamard construction transfers to Bianchi $\mathrm{VI_0}$:
\begin{itemize}
  \item \textbf{For the massive conformally coupled scalar}: in principle, conditional on
        Lemmas C and D.  The overall strategy is sound and the Sol-specific difficulties
        (Bessel fiber, Mathieu eigenvalue) are technically hard but not fundamentally
        obstructed.
  \item \textbf{For the massless scalar}: genuinely obstructed at IR (Obstruction A).
        This is not a Sol-specific failure but a universal feature of massless SLE
        constructions (already in BN23 and Olbermann).
\end{itemize}

\subsection{Remaining technical work}

A realistic breakdown of the open work:

\begin{center}
\begin{tabular}{lll}
\hline
Item & Nature & Estimated effort \\
\hline
Lemma C (Mathieu eigenvalue positivity) & Spectral theory & 2--4 weeks \\
Lemma D (Adiabatic order for Bessel-$K$) & Asymptotic analysis & 4--8 weeks \\
Obstruction E (WF set full proof) & Microlocal analysis & 3--6 weeks \\
Obstruction F (Quotient / compactification) & Harmonic analysis & 2--4 weeks \\
IR cutoff / mass: full theorem statement & Functional analysis & 1--2 weeks \\
\hline
Total (sequential, one author) & & $\sim 3$--6 months \\
\hline
\end{tabular}
\end{center}

\subsection{Honest assessment}

The classification of this piste as ``the hardest among Type A'' in ECI v6.0.24--26 is
justified.  The Sol group is the first Bianchi case where:
(i)~plane-wave Fourier analysis is replaced by a two-parameter family of
infinite-dimensional representations with fiber-direction Bessel functions;
(ii)~the temporal mode equation depends on an eigenvalue that is itself a solution of
a Mathieu equation;
(iii)~the compactification involves an Anosov lattice whose spectral theory is
substantially harder than the flat torus of Bianchi I.

Under favourable conditions (single author, specialist in microlocal analysis and
harmonic analysis on solvable groups), a complete paper could be written in
$\sim 4$--$6$ months.  Under realistic conditions (need to learn Mathieu/Bessel
asymptotics, BN23 framework, and Sol harmonic analysis simultaneously), a 9--12 month
timeline is more appropriate.

\paragraph{Recommendation.}
Proceed with the massive scalar case.  Write Lemma C first (spectral theory, computationally
accessible via Mathieu equation references).  Lemma D (adiabatic order with Bessel-$K$)
is the key novel contribution and will determine whether this becomes a standalone paper
or part of a larger series.

%% ---------------------------------------------------------------
\section*{References}
\addcontentsline{toc}{section}{References}
%% ---------------------------------------------------------------

\begin{thebibliography}{99}

\bibitem{AK71}
  L.~Auslander and B.~Kostant,
  ``Polarization and unitary representations of solvable Lie groups,''
  \emph{Inventiones Mathematicae} \textbf{14} (1971), 255--354.
  \textit{[arXiv:not available; Springer DOI: 10.1007/BF01389744.
   Reference VERIFIED via Springer Nature database.]}

\bibitem{AV13}
  Zh.~Avetisyan and R.~Verch,
  ``Explicit harmonic and spectral analysis in Bianchi I-VII type cosmologies,''
  \emph{Classical and Quantum Gravity} \textbf{30} (2013), 155006.
  arXiv:1212.6180 [math-ph].
  \textit{[arXiv ID VERIFIED.  Covers Bianchi I--VII including VI$_0$.]}

\bibitem{BN23}
  R.~Banerjee and M.~Niedermaier,
  ``States of Low Energy on Bianchi I spacetimes,''
  \emph{Journal of Mathematical Physics} \textbf{64} (2023), 113503.
  arXiv:2305.11388 [math-ph].
  \textit{[arXiv ID VERIFIED.  Main reference for SLE on anisotropic backgrounds.]}

\bibitem{BianchiClass}
  E.~Bianchi,
  ``Sugli spazi a tre dimensioni che ammettono un gruppo continuo di movimenti
  (On the spaces of three dimensions that admit a continuous group of movements),''
  \emph{Memorie di Matematica e di Fisica della Societa Italiana delle Scienze} (1898).
  Modern treatment: R.~T.~Jantzen,
  ``The Bianchi classification in the Sch\"ucking--Behr approach,''
  in \emph{Golden Oldie} reprints (GRG journal).

\bibitem{Kir04}
  A.~A.~Kirillov,
  \emph{Lectures on the Orbit Method},
  Graduate Studies in Mathematics \textbf{64},
  American Mathematical Society, Providence, RI, 2004.
  ISBN: 978-0-8218-3530-2.
  \textit{[Reference VERIFIED: AMS GSM vol.~64, 408 pp.]}

\bibitem{Olb07}
  H.~Olbermann,
  ``States of low energy on Robertson-Walker spacetimes,''
  \emph{Classical and Quantum Gravity} \textbf{24} (2007), 5011--5030.
  arXiv:0704.2986 [gr-qc].
  \textit{[arXiv ID VERIFIED.  Original SLE definition.]}

\bibitem{Rad96}
  M.~J.~Radzikowski,
  ``Micro-local approach to the Hadamard condition in quantum field theory on
  curved space-time,''
  \emph{Communications in Mathematical Physics} \textbf{179} (1996), 529--553.
  \textit{[DOI: 10.1007/BF02100096.  Reference VERIFIED.]}

\end{thebibliography}

\end{document}
